\smallskip
\noindent\emph{The weakest one-bit class (resolving the open piece of
Conjecture~\ref{conj:gen}).}
Theorem~\ref{thm:conj1-dichotomy} fixes the supplement at \emph{one bit}; what remained open is
the \emph{weakest falsifiable class} on which that bit suffices. We settle it conditionally. On
$D$ write $a(x):=2\eta_a(x)-1$ with $\eta_a(x)=\Pr_T(f_a{=}Y\mid x)$, so $\Delta=\int_D a\,d\mu$
(Lemma~\ref{lem:reduction}). Let $m(x):=2s(x)-1$ be the observable relative (calibrated-score)
margin, so $M=\tfrac12\E_{\mu\mid D}[m]$, and let the observable \emph{relative-calibration
field} be $c(x):=w(x)\,(m(x)-m^\star)$ with observable nonnegative weight $w(x)$ and observable
neutral anchor $m^\star$ (the margin at which $s=\tfrac12$). Put $T_E:=\int_D c\,d\mu$
(observable), and let $E$ denote the law of all label-free observables; two targets are
\emph{evidence-identical} when they share $E$ (equivalently $(\mu,m,s)$, since predictions and
$s$ are fixed maps of $x$).

\begin{definition}[Margin-monotone relative-calibration class]\label{def:cmono}
$\mathcal C_{\mathrm{mono}}$ is the class of targets for which $a(x)=\sigma\,c(x)$ $\mu$-a.e.\ on
$D$ for a single global orientation $\sigma\in\{\pm1\}$; equivalently the benefit is
single-crossing in the observable margin at $m^\star$ ($\sign a(x)=\sigma\,\sign(m(x)-m^\star)$)
with calibrated magnitude $|a(x)|=|c(x)|$. The class is \emph{falsifiable but not verifiable}: a
finite labelled probe can reject ``$a=\pm c$'', yet $\sigma$ is never label-free identified.
\end{definition}

\begin{assumption}[General position]\label{asm:genpos}
Within an evidence fibre the regions of $D$ whose benefit sign is free carry comparable,
not-fully-$E$-pinned benefit mass: $T_E\neq0$ and no proper subcollection of free regions has an
$E$-certifiable dominant total $\int|a|\,d\mu$.
\end{assumption}

\begin{lemma}[Structural sign-freedom $=$ bit complexity]\label{lem:bitcomplexity}
Fix $E$ and a falsifiable class $\mathcal C$, and let $K(\mathcal C,E)$ be the number of disjoint
$\mu$-positive regions of $D$ whose benefit signs are jointly free in $\mathcal C$ on the fibre
$E$. Under Assumption~\ref{asm:genpos}, a falsifiable structural supplement of $b$ bits pins
$\sign\Delta$ across the fibre \emph{iff} $b\ge K(\mathcal C,E)$.
\end{lemma}
\begin{proof}
If $b\ge K$, declaring the orientation of each free region fixes the sign field; the rest of $a$
being $E$-pinned, $\Delta=\int_D a\,d\mu$ then has $E$-computable sign. If $b<K$, some free region
$R$ is undeclared; the two targets agreeing off $R$ and differing in $\sign a$ on $R$ are
evidence-identical and share all $b$ declared bits, while by comparability of $\int_R|a|$ with the
complement (Assumption~\ref{asm:genpos}) their benefits have opposite sign---no decoder of the $b$
declared bits separates them.
\end{proof}

\begin{theorem}[Conditional resolution of Conjecture~\ref{conj:gen}: a weakest one-bit class]\label{thm:cmono-weakest}
\textnormal{(i) Sufficiency \emph{(unconditional)}:} every target in $\mathcal C_{\mathrm{mono}}$ has
$\sign\Delta=\sigma\,\sign(T_E)$ with $T_E$ observable, so $E$ and the single orientation bit
$\sigma$ determine $\sign\Delta$.
\textnormal{(ii) Necessity \emph{(unconditional)}:} $\sigma$ and $-\sigma$ are evidence-identical
($\Law(E\mid\sigma)=\Law(E\mid-\sigma)$) yet give $\Delta$ and $-\Delta$, so by Le~Cam at least
one bit is required. Thus \emph{exactly one bit certifies $\sign\Delta$ on
$\mathcal C_{\mathrm{mono}}$, with no further assumption}.
\textnormal{(iii) Minimality \emph{(under Assumption~\ref{asm:genpos})}:} every falsifiable
$\mathcal C\supsetneq\mathcal C_{\mathrm{mono}}$ contains a non-margin-monotone target with
$K(\mathcal C,E)\ge2$, so by Lemma~\ref{lem:bitcomplexity} at least two bits are required. Hence
$\mathcal C_{\mathrm{mono}}$ is \emph{a} weakest falsifiable one-bit class---minimal in the
single-crossing, calibrated-magnitude direction; it is not unique (Remark~\ref{rmk:genpos}).
\end{theorem}
\begin{proof}
(i) $\Delta=\int_D a\,d\mu=\sigma\int_D c\,d\mu=\sigma T_E$. (ii) Sending $\sigma\mapsto-\sigma$ is
the $D$-local label complement of Theorem~\ref{thm:conj1-dichotomy}: it fixes $(\mu,m,s)$, hence
$E$, while negating $a$ and thus $\Delta$; the two-point bound of Theorem~\ref{thm:imp} gives
minimax committal error $\tfrac12$. (iii) A target outside $\mathcal C_{\mathrm{mono}}$ violates
$\sign a=\sigma\,\sign(m-m^\star)$ on a $\mu$-positive region, so its sign field is not one global
orientation; in the swap-closed (falsifiable) fibre at least two regions are independently
signable, i.e.\ $K\ge2$, and Lemma~\ref{lem:bitcomplexity} applies.
\end{proof}

\begin{proposition}[Explicit non-identifiability witness]\label{prop:cmono-witness}
For \emph{each} candidate single structural bit there are two targets with identical evidence law
($\TV=0$) sharing that bit's value but with opposite $\sign\Delta$. Take $D=R_1\cup R_2$,
$\mu(R_1)=\mu(R_2)=\tfrac12$, with $R_1$ calibrated ($c_1{=}1$) and $R_2$ at the anchor
($c_2{=}0$, no observable benefit signal). For $\beta=\sign a_1$: $a=(+1,-2)$ vs.\ $(+1,+2)$ give
$\Delta=-\tfrac12$ vs.\ $+\tfrac32$. For $\beta=\sign a_2$: $a=(+1,+\tfrac15)$ vs.\
$(-1,+\tfrac15)$ give $\Delta=+\tfrac35$ vs.\ $-\tfrac25$. No single bit certifies $\sign\Delta$,
so two bits are necessary off $\mathcal C_{\mathrm{mono}}$.
\end{proposition}

\begin{remark}[What is unconditional, and what rests on general position]\label{rmk:genpos}
Parts (i)--(ii)---that one declared bit certifies $\sign\Delta$ on $\mathcal C_{\mathrm{mono}}$---hold
\emph{unconditionally}. Only minimality (iii) uses Assumption~\ref{asm:genpos}: if it fails (one
region's benefit mass is $E$-certifiably dominant) that region's orientation alone certifies
$\sign\Delta$ even outside $\mathcal C_{\mathrm{mono}}$, moving the boundary. Equivalently,
Assumption~\ref{asm:genpos} is the requirement that the certificate be \emph{domination-free}
(minimax-sound over the unobserved magnitudes); so the General-Position statement is only
\emph{conditional}; the unconditional weakest classes are characterized in
Theorem~\ref{thm:uncond-weakest} as an explicit finite family of dominance polytopes. Even under
Assumption~\ref{asm:genpos}, $\mathcal C_{\mathrm{mono}}$ is only \emph{a} weakest one-bit class:
any calibrated sign-aligned field $\{a=\sigma h:\,h\ \text{observable}\}$ is equally one-bit and
is incomparable to $\mathcal C_{\mathrm{mono}}$. All six checks are machine-checked by a
deterministic release validator: the unconditional core (sufficiency/necessity
$20000/20000$, $|\Delta-\sigma T_E|=0$ exactly), the explicit minimality witnesses of~(iii), the
general-position ablation, and the non-uniqueness check.

\smallskip
\noindent\emph{Why general position cannot simply be dropped (the precise obstruction, and
exactly what is open).} The minimal counterexample is two equal-mass regions
$D=R_0\cup R_1$, $\mu(R_0)=\mu(R_1)=\tfrac12$, with the observable field $c=+1$ on the
\emph{calibrated} region $R_0$ and $c=0$ on the \emph{anchor} region $R_1$ (no observable benefit
signal there, so $R_1$ is sign-free in any evidence fibre). Enlarge $\mathcal C_{\mathrm{mono}}$ to
the falsifiable class $\mathcal C_{\mathrm{dom}}=\{a_{R_0}=\sigma,\ |a_{R_1}|\le\rho\}$ for a
declared constant $\rho<1$. Since $\Delta=\tfrac12(a_{R_0}+a_{R_1})$ and $|a_{R_1}|\le\rho<1=|a_{R_0}|$,
the calibrated region dominates and $\sign\Delta=\sigma$ for \emph{every} member, so one declared
bit certifies $\sign\Delta$ on all of $\mathcal C_{\mathrm{dom}}$. Yet $\mathcal C_{\mathrm{dom}}
\supsetneq\mathcal C_{\mathrm{mono}}$ (it contains non-margin-monotone targets with $a_{R_1}\neq0$,
e.g.\ $a=(+1,\pm\tfrac12)$ giving $\Delta=\tfrac34,\tfrac14$, both positive) and is falsifiable (a
labelled probe rejects targets that are uncalibrated on $R_0$ or violate $|a_{R_1}|\le\rho$). This
\emph{contradicts} minimality once an $E$-certifiable dominant region is permitted: dropping
$a=(+1,\pm2)\Rightarrow\Delta=+\tfrac32,-\tfrac12$ shows that the moment $\rho\ge1$ the anchor
region can flip $\sign\Delta$ and two bits become necessary ($K{=}2$). The obstruction to an
\emph{unconditional} weakest class is therefore exactly this: characterising minimality across the
\emph{entire} lattice of dominant-region thresholds $\rho$ (where the weakest one-bit class is not
$\mathcal C_{\mathrm{mono}}$ but a $\rho$-dependent family), which Assumption~\ref{asm:genpos}
sidesteps by fiat. Theorem~\ref{thm:uncond-weakest} (below) resolves it: the $\rho$-lattice is the
slice $\{r_0{=}1,\,r_1{=}\rho\}$ of the single dominance polytope $W^\star{=}\{r_0\ge r_1\}$, with
$\rho{=}1$ its boundary. The construction is machine-checked
(\texttt{val\_conj1\_genpos.py}, seeds fixed): $\mathcal C_{\mathrm{dom}}$ strictly contains
$\mathcal C_{\mathrm{mono}}$, is falsifiable (reject rate $1.0$), is one-bit-certified (error rate
$0$ over $2\times10^5$ targets), and loses one-bit certifiability without the domination bound
(error rate $\approx\tfrac14$).
\end{remark}

\smallskip
\noindent\emph{Unconditional resolution (General Position removed).}
Dropping Assumption~\ref{asm:genpos} reveals, for each orientation pattern, a single maximal one-bit
class---a \emph{dominance polytope}---of which the $\rho$-family of Remark~\ref{rmk:genpos} is one
slice.

\begin{definition}[Orientation pattern and canonical class]\label{def:orient}
On a fibre $E$ an \emph{orientation pattern} is a tied set $P\subseteq\{i:c_i\neq0\}$ with a tied
sign-pattern $s\in\{\pm\}^{P}$; the free set is $Z=D\setminus P$. For an observable magnitude region
$W\subseteq[0,1]^{|D|}$ the \emph{canonical class} is
$\mathcal C(P,s,W)=\{a:\ \sign a_i=\sigma\,s_i\sign c_i\ (i\in P),\ (|a_j|)_j\in W\}$ with one
declared bit $\sigma$. Write $T(r)=\sum_{i\in P}s_i\sign(c_i)\mu_i r_i$ and
$G(r)=\sum_{i\in Z}\mu_i r_i$.
\end{definition}

\begin{lemma}[Canonical form of a falsifiable class]\label{lem:canonform}
A class is falsifiable \emph{iff} it equals some $\mathcal C(P,s,W)$; anchors ($c_i=0$) lie
necessarily in $Z$. \textnormal{(Swap-closure, Thm.~\ref{thm:conj1-dichotomy}(iii), forces the
constraint on $Z$ to depend only on $|a|$; forcing an anchor sign is not evidence-definable. Both
aligned ($s_i{=}+$) and anti-aligned ($s_i{=}-$) ties are falsifiable since $\sigma$ is
unidentified.)}
\end{lemma}

\begin{theorem}[Unconditional one-bit criterion; weakest classes]\label{thm:uncond-weakest}
Let $\mathcal C(P,s,W)$ be falsifiable. \textnormal{(i) (Dominance margin.)} One declared bit
certifies $\sign\Delta$ on $\mathcal C(P,s,W)$ \emph{iff}
$\inf_{r\in W}\!\big(T(r)-G(r)\big)\ge0$ (with some member $\Delta>0$) or
$\sup_{r\in W}\!\big(T(r)+G(r)\big)\le0$. \textnormal{(ii) (Weakest class.)} For a fixed pattern
$(P,s)$ the unique maximal one-bit class is the dominance polytope
$\mathcal C^\star=\mathcal C(P,s,W^\star)$, $W^\star=\{r\in[0,1]^{|D|}:T(r)\ge G(r)\}$;
$\mathcal C_{\mathrm{mono}}$, every $\mathcal C_{\mathrm{dom}}(\rho\le1)$, and every box one-bit
class are proper subclasses of one $\mathcal C^\star$. \textnormal{(iii) (Finite family.)} The set
of weakest falsifiable one-bit classes is the explicit finite family $\{\mathcal C^\star\}$ indexed
by $(P\subseteq\{c\neq0\},\,s\in\{\pm\}^{P})$ modulo the global flip $\sigma\mapsto-\sigma$ (at most
$3^{|\{c\neq0\}|}$ patterns); distinct patterns give pairwise incomparable maxima, so there is no
\emph{unique} weakest class. $\mathcal C_{\mathrm{mono}}$ is the all-tied all-aligned member, and
Assumption~\ref{asm:genpos} is exactly the face on which the family collapses to
$\mathcal C_{\mathrm{mono}}$, recovering Thm.~\ref{thm:cmono-weakest}(iii).
\end{theorem}
\begin{proof}
$\Delta$ is linear, so over $\mathcal C(P,s,W)$ at bit $\sigma=+1$ it ranges over
$[\inf_W(T-G),\sup_W(T+G)]$ (free signs independent by swap-closure; the $\sigma=-1$ fibre is its
negation); a single bit certifies iff each fibre's strict sign set is a singleton, giving (i). For
$W^\star$, $\inf_{W^\star}(T-G)\ge0$ by construction, and any $r_0$ with $T(r_0)<G(r_0)$ admits,
with opposing free signs, a member with $\Delta<0$ while $W^\star$ admits $\Delta>0$, so $W^\star$
is maximal; $\mathcal C_{\mathrm{dom}}(\rho)$ sits in $W^\star$ iff $\rho\le1$, with $\rho=1$ the
boundary, giving (ii). Incomparability (iii) is witnessed on $c=(1,1,0)$ by $a=(0,\tfrac13,0)$ and
$a=(1,-\tfrac13,\tfrac13)$. On the general-position face the free block can overpower the tied block
unless $Z$-magnitudes vanish, leaving only $\mathcal C_{\mathrm{mono}}$.
\end{proof}

\noindent\emph{Verification.} The criterion~(i) was checked against exact brute-force vertex truth
on $2.8{\times}10^{5}$ box fibres and $3.2{\times}10^{3}$ non-box integer polytopes (exact-LP
$\inf/\sup$) with $0$ mismatches; sufficiency on $W^\star$ over $1.2{\times}10^{5}$ members
($0$ errors), necessity just outside $W^\star$ (error rate $\approx\tfrac12$), and the
$\mathcal C_{\mathrm{dom}}(\rho)$ collapse ($\rho{\le}1$ one-bit, $\rho{>}1$ lost) are
machine-checked (\texttt{val\_unconditional\_weakest.py}, seeds fixed; fails loudly). The single
non-tautological input is Lemma~\ref{lem:canonform}, which rests on the swap involution
(Thm.~\ref{thm:conj1-dichotomy}(iii)).
